\section{Introduction}

For a given portfolio, we can estimate its expectation of return. However, the actual return of an investment in this portfolio may be larger or smaller than this expectation.
Therefore, the dispersion of outcomes is also a vital information for an investor.
A smaller dispersion value may translate in lower probability for outcomes far below the expected return.

A risk-based optimal portfolio is a portfolio that, in some sense, has the best combination between the expected return and the dispersion of outcomes.

In the context of portfolio optimizations, the dispersion of outcomes is often called risk.\footnote{Do not confuse it with the risk
notion in the context of risk management, where risk is understood as the amount of capital reserves needed to cover the potential losses of an investment,
up to a certain level of confidence. These two notions are closely related to each other but different.}
The quantity that measures the dispersion of outcomes is called dispersion or risk measure.\footnote{Some authors prefer to call
the dispersion measure as deviation. However, we all refer to the same mathematical object.}
Where there are no confusions, we will use both terms interchangeable.

It was Markowitz in 1952 \cite{Markowitz} who pointed out that for a given expected rate of return
there is a portfolio that has the smallest value of variance.
He called this portfolio {\it efficient}. Moreover, for efficient portfolios the variance
is increasing with the value of the expected rate of return.
Therefore, an investment is subject to a tradeoff between the size of the expected payoff and the risk that is exposed to.
In Markowitz's work, variance was playing the role of
dispersion measure. And thus, the {\it Modern Portfolio Theory} was born. These days the variance based optimal strategies
are often called mean-variance (\MV) strategies. They will be discussed in Sec.~\ref{S:MV}.

Since 1952, many risk-based portfolio strategies were formulated simply by replacing the variance with other
dispersion measures; sometimes more layers of mathematical refinements were added.

In this paper the risk-based portfolio strategies are categorized by the families of underlining dispersion measures.
We distinguish 5 such families: tail, downside, below threshold, inter-observations distance and central moment dispersion measures.
In each case we seek the most general form in terms of $L^p$-norms.
Numerical algorithms are presented for $p=1$ and $p=2$. In these cases, the portfolio optimizations are
formulated as numerically efficient LP (Linear Programming), QP (Quadratic Programming)
and SOCP (Second Order Cone Programming) problems.


Section~\ref{strategies} presents the generic portfolio optimization strategies.
We had identified the following strategies:
\begin{itemize}
  \item minimization of risk for targeted expected rate of return value,
  \item minimum risk portfolio,
  \item maximization of expected rate of return for targeted risk value,
  \item maximization of expected rate of return for fixed risk-aversion factor,
  \item maximization of generalized Sharpe ratio,
  \item minimization of inverse generalized Sharpe ratio.
\end{itemize}
Although not all of them are mutually exclusive, they lead to different numerical algorithms.
In the next sections, these generic strategies will be applied to the following dispersion measures:
\begin{itemize}
  \item mixture Conditional Value at Risk (\mCVaR),  Sec.~\ref{S:mCVaR},
  \item mixture Second Moment Coherent Risk (\mSMCR),  Sec.~\ref{S:mSMCR},
  \item m-level Mean Absolute Deviation, (\mMAD),  Sec.~\ref{S:mMAD},
  \item m-level Lower Semi-Deviation (\mLSD),  Sec.~\ref{S:mLSD},
  \item mixture Below Target Absolute Deviation (\mBTAD),  Sec.~\ref{S:mBTAD},
  \item mixture Below Target Semi-Deviation (\mBTSD),  Sec.~\ref{S:mBTSD},
  \item Gini index (\GINI),  Sec.~\ref{S:GINI},
  \item standard deviation (\SD),  Sec.~\ref{S:SD},
  \item variance (\VAR) - leading to Mean-Variance model, \MV,  Sec.~\ref{S:MV},
\end{itemize}
resulting in concrete numerical algorithms.

Note that, some of the algorithms have well recognized popular names. An example is the
Omega and Sortino optimal portfolios which in our classification are special cases
of \mBTAD-Sharpe and \mBTSD-Sharpe optimal portfolios, respectively.
In all cases we have adopted the popular names with the appropriate caveat regarding generalizations.
Most of the algorithms have been previously mentioned in the literature, in one way or the other.
However, for some of them we were not able to identify relevant literature. Given our systematic classification,
we were able to fill in the gaps. For example, the  \mLSD\ based optimal portfolio strategies appear to be presented here for the first time,
although they are clear extensions of \mMAD, where $L^1$ is replaced by $L^2$-norm and LP by SOCP problems.
A complete list of all the algorithms is in the table of contents.

We made available the implementation of these algorithms in an open-source Python package called \azapy.
The package can be installed in a Python environment simple by calling {\bf pip install azapy}.
The technical documentation of the library can be found at \azapydoc\ and its source code at \azapycode.

Simple Python scripts illustrating the use of \azapy\ library are accompanying the portfolio optimization
strategies as well as in-sample and out-of-sample portfolio simulations.
A brief introduction to \azapy\ package is presented in Appendix~\ref{AZAPY}.
Technical details about the portfolio backtesting (\ie\ out-of-sample testing) methodology are presented in Appendix~\ref{backtesting}.

Regarding our terminology, the time horizon for a dispersion measure is simple defined as the duration of the rate of return argument (\ie\
a monthly rate produces a dispersion measure with a time horizon of one month,
etc.).
The calibration period for a dispersion measure is the look-back period of historical observations used to evaluate the dispersion.
Usually, it is measured in years. Often, in this paper are mentioned dispersion measures with time horizon of one quarter and
calibration period of 3.25 years.

Technical sections, presenting various optimization algorithms, are meant to be self-sufficient, with a minimum number of cross references to other parts of the paper. To facilitate the navigation of this document, at the end of each section there is a hyperlink pointing to the table of contents, from where one can jump to other sections.

\BackToContents



